\documentclass[diploma]{softlab-thesis}

%%%
%%%  Add and configure the packages that you need for your thesis
%%%

\usepackage{minted}

% Hyperref needs to be the last package to use!
\PassOptionsToPackage{hyphens}{url}
\usepackage[xetex,%
    pdfpagemode=UseOutlines,pdfstartview=FitW,%
    colorlinks=true,linkcolor=blue,citecolor=blue,urlcolor=blue,%
    bookmarks=true,bookmarksopen=true,bookmarksnumbered=false,%
    pdfencoding=auto,unicode=true,hyperfootnotes=true,%
    hypertexnames=false,%
  ]{hyperref}


%%%
%%%  The document
%%%

\begin{document}

%%%  Title page

\frontmatter

\title{Μοντελοποίηση Αγορών Ηλεκτρικής Ενέργειας με Julia/JuMP}
\author{Ιωάννης Κ. Γεωργακόπουλος}
\authoren{Ioannis K. Georgakopoulos}
\date{Φεβρουάριος 2025}
\datedefense{28}{2}{2025} % Famous Last Words

\supervisor{Άρης-Ευάγγελος Δημέας}
\supervisorpos{Επίκουρος Καθηγητής Ε.Μ.Π.}

\committeeone{Άρης-Ευάγγελος Δημέας}
\committeeonepos{Επίκουρος Καθηγητής Ε.Μ.Π.}
\committeetwo{Γεώργιος Κορρές}
\committeetwopos{Καθηγητής Ε.Μ.Π.}
\committeethree{Γεώργιος Νικολάου}
\committeethreepos{Αν. Καθηγητής Ε.Κ.Π.Α.}

\TRnumber{CSD-SW-TR-42-17}  % number-year, ask nickie for the number
\department{Τομέας Ηλεκτρικής Ισχύος}

\maketitle


%%%  Abstract, in Greek

\begin{abstractgr}%
  Σκοπός της παρούσας εργασίας είναι αφενός η σχεδίαση μίας απλής
  γλώσσας υψηλού επιπέδου με υποστήριξη για προγραμματισμό με
  αποδείξεις, αφετέρου η υλοποίηση ενός μεταγλωττιστή για τη γλώσσα
  αυτή που θα παράγει κώδικα για μία γλώσσα ενδιάμεσου επιπέδου
  κατάλληλη για δημιουργία πιστοποιημένων εκτελέσιμων.

  Στη σημερινή εποχή, η ανάγκη για αξιόπιστο και πιστοποιημένα ασφαλή
  κώδικα γίνεται διαρκώς ευρύτερα αντιληπτή. Τόσο κατά το παρελθόν όσο
  και πρόσφατα έχουν γίνει γνωστά προβλήματα ασφάλειας και
  συμβατότητας προγραμμάτων που είχαν ως αποτέλεσμα προβλήματα στην
  λειτουργία μεγάλων συστημάτων και συνεπώς οικονομικές επιπτώσεις
  στους οργανισμούς που τα χρησιμοποιούσαν. Τα προβλήματα αυτά
  οφείλονται σε μεγάλο βαθμό στην έλλειψη δυνατότητας προδιαγραφής και
  απόδειξης της ορθότητας των προγραμμάτων που χαρακτηρίζει τις
  σύγχρονες γλώσσες προγραμματισμού. Για το σκοπό αυτό, έχουν προταθεί
  συστήματα πιστοποιημένων εκτελέσιμων, στα οποία έχουμε τη δυνατότητα
  να προδιαγράφουμε την ορθότητα των προγραμμάτων, και να παρέχουμε
  μία τυπική απόδειξη αυτής, η οποία μπορεί να ελεγχθεί μηχανιστικά
  πριν το χρόνο εκτέλεσης.

  Τα συστήματα που έχουν προταθεί είναι ενδιάμεσου επιπέδου οπότε η
  διαδικασία προγραμματισμού σε αυτά είναι ιδιαίτερα πολύπλοκη. Οι
  γλώσσες υψηλού επιπέδου που συνοδεύουν αυτά τα συστήματα, ενώ είναι
  ιδιαίτερα εκφραστικές, παραμένουν δύσκολες στον προγραμματισμό.  Μία
  απλούστερη γλώσσα υψηλού επιπέδου, όπως αυτή που προτείνουμε σε αυτή
  την εργασία, θα επέτρεπε ευρύτερη εξάπλωση του συγκεκριμένου
  ιδιώματος προγραμματισμού.

  Στη γλώσσα που προτείνουμε, ο προγραμματιστής προδιαγράφει τη μερική
  ορθότητα του προγράμματος, δίνοντας προσυνθήκες και μετασυνθήκες για
  τις παραμέτρους και τα αποτελέσματα των συναρτήσεων που ορίζει.
  Επίσης δίνει ένα σύνολο θεωρημάτων βάσει του οποίου κατασκευάζονται
  αποδείξεις της ορθής υλοποίησης και χρήσης των συναρτήσεων αυτών. Ως
  μέρος της εργασίας, έχουμε υλοποιήσει σε γλώσσα OCaml ένα
  μεταφραστή αυτής της γλώσσας στο σύστημα πιστοποιημένων
  εκτελέσιμων NFLINT.

  Επιτύχαμε να διατηρήσουμε τη γλώσσα κοντά στο ύφος των ευρέως
  διαδεδομένων συναρτησιακών γλωσσών, καθώς και να διαχωρίσουμε τη
  φάση προγραμματισμού από τη φάση απόδειξης της ορθότητας των
  προγραμμάτων. Έτσι ένας μέσος προγραμματιστής μπορεί εύκολα να
  προγραμματίζει στη γλώσσα που προτείνουμε με τον τρόπο που ήδη
  γνωρίζει, και ένας γνώστης μαθηματικής λογικής να αποδεικνύει σε
  επόμενη φάση την μερική ορθότητα των προγραμμάτων. Ως απόδειξη της
  πρακτικότητας της προσέγγισης αυτής, παραθέτουμε ένα σύνολο
  παραδειγμάτων στη γλώσσα με απόδειξη μερικής ορθότητας.
\begin{keywordsgr}
  Γλώσσες προγραμματισμού,
  Προγραμματισμός με αποδείξεις,
  Ασφαλείς γλώσσες προγραμματισμού,
  Πιστοποιημένος κώδικας.
\end{keywordsgr}
\end{abstractgr}


%%%  Abstract, in English

\begin{abstracten}%
  The purpose of this diploma dissertation is on one hand the design
  of a simple high-level language that supports programming with
  proofs, and on the other hand the implementation of a compiler for
  this language. This compiler will produce code for an
  intermediate-level language suitable for creating certified
  binaries.

  The need for reliable and certifiably secure code is even more
  pressing today than it was in the past. In many cases, security and
  software compatibility issues put in danger the operation of large
  systems, with substantial financial consequences. The lack of a
  formal way of specifying and proving the correctness of programs that
  characterizes current programming languages is one of the main reasons
  why these issues exist. In order to address this problem, a number of
  frameworks with support for certified binaries have recently been
  proposed. These frameworks offer the possibility of specifying and
  providing a formal proof of the correctness of programs. Such a proof
  can easily be checked for validity before running the program.

  The frameworks that have been proposed are intermediate-level in
  nature, thus the process of programming in these is rather cumbersome.
  The high-level languages that accompany some of these frameworks,
  while very expressive, are hard to use. A simpler high-level language,
  like the one proposed in this dissertation, would enable further use
  of this programming idiom.

  In the language we propose, the programmer specifies the partial
  correctness of a program by annotating function definitions with pre-
  and post-conditions that must hold for their parameters and results.
  The programmer also provides a set of theorems, based on which proofs
  of the proper implementation and use of the functions are constructed.
  An implementation in OCaml of a compiler from this language to the
  NFLINT certified binaries framework was also completed as part of this
  dissertation.

  We managed to keep the language close to the feel of the current
  widespread functional languages, and also to fully separate the
  programming stage from the correctness-proving stage. Thus an average
  programmer can program in a familiar way in our language, and later an
  expert on formal logic can prove the semi-correctness of a program.
  As evidence of the practicality of our design, we provide a number of
  examples in our language with full semi-correctness proofs.
\begin{keywordsen}
  Programming languages,
  Programming with proofs,
  Secure programming languages,
  Certified code.
\end{keywordsen}
\end{abstracten}


%%%  Acknowledgements

\begin{acknowledgementsgr}
  Ευχαριστώ θερμά τον επιβλέποντα καθηγητή αυτής της διατριβής,
  κ.~Νίκο Παπασπύρου, για τη συνεχή καθοδήγηση και εμπιστοσύνη
  του. Ευχαριστώ επίσης τα μέλη της συμβουλευτικής επιτροπής,
  κ.κ.~Νίκο Παπαδόπουλο και Γιώργο Νικολάου για την πρόθυμη και
  πάντα αποτελεσματική βοήθειά τους, τις πολύτιμες συμβουλές και
  τις χρήσιμες συζητήσεις που είχαμε.  Θέλω να ευχαριστήσω ακόμα
  τον συμφοιτητή και φίλο Πέτρο Πετρόπουλο, ο οποίος με βοήθησε σε
  διάφορα στάδια αυτής της εργασίας.  Θα ήθελα τέλος να ευχαριστήσω
  την οικογένειά μου και κυρίως τους γονείς μου, οι οποίοι με
  υποστήριξαν και έκαναν δυνατή την απερίσπαστη ενασχόλησή μου τόσο
  με την εκπόνηση της διπλωματικής μου, όσο και συνολικά με τις
  σπουδές μου, καθώς και το Νίκο Καζαντζάκη που αποτέλεσε για μένα
  μοναδική πηγή έμπνευσης, όταν δεν ήξερα τι να γράψω στην εισαγωγή
  της εργασίας μου.
\end{acknowledgementsgr}

\begin{acknowledgementsen}
  I would like to thank all the people who supported my work and helped me get
  results of better quality.  I am also grateful to the supervisor of my thesis
  and the members of my committee for their patience and support in overcoming
  numerous obstacles I have been facing through my research.  I would also like
  to thank my fellow students for their feedback, cooperation and of course
  friendship.

  Last but not least, I would like to thank my family, my parents and my
  sister, for supporting me spiritually throughout writing this thesis and
  my life in general, and Nikos Kazantzakis for providing the inspiration
  that I needed when I didn't know what to write in my introduction.
\end{acknowledgementsen}


%%%  Various tables

\tableofcontents
%\listoftables
\listoffigures
%\listofalgorithms


%%%  Main part of the book

\mainmatter

\chapter{Εισαγωγή}

Εδώ εισαγωγή


\section{Σκοπός της εργασίας}

Εδώ σκοπός

\section{Δέσμευση Μονάδων}

Για να κατανοήσουμε τη φιλοσοφία του ευρωπαϊκού μοντέλου και του αλγορίθμου EUPHEMIA, πρέπει πρώτα να
εξετάσουμε το πρόβλημα δέσμευσης μονάδων. Το πρόβλημα δέσμευσης μονάδων (Unit Commitment Problem)
στα Συστήματα Ηλεκτρικής Ενέργειας είναι το πρόβλημα προγραμματισμού της κατάστασης λειτουργίας των 
μονάδων παραγωγής, το οποίο περιλαμβάνει τους χρονικά εξαρτημένους περιορισμούς, 
καθώς και τα χρονικά εξαρτημένα κόστη.
Η εξισορρόπηση της παραγωγής και της κατανάλωσης αποτελεί μία συνεχή προσπάθεια για την ορθή και 
αξιόπιστη λειτουργία των Συστημάτων Ηλεκτρικής Ενέργειας. Δεδομένων των λειτουργικών περιορισμών 
των μονάδων (περιορισμοί ράμπας, ελάχιστοι χρόνοι εντός και λειτουργίας, κ.α.), η λήψη ορισμένων 
αποφάσεων πρέπει να πραγματοποιηθεί πριν τον πραγματικό χρόνο, για τη διασφάλιση της εξισορρόπησης 
της παραγωγής με την κατανάλωση κάθε χρονική στιγμή. 
Για τον σκοπό αυτό, είναι ευρέως διαδεδομένη η διάκριση μεταξύ προγραμματισμού επόμενης μέρας και 
κατανομής σε πραγματικό χρόνο. Ο προγραμματισμός επόμενης μέρας βασίζεται σε μία πρόγνωση της ζήτησης 
και ορίζει το πρόγραμμα αργών μονάδων, με χαρακτηριστικό παράδειγμα τις λιγνιτικές, ακολουθώντας το 
μοντέλο δέσμευσης μονάδων.  Στη φάση της κατανομής πραγματικού χρόνου, ο εκάστοτε διαχειριστής 
εξισορροπεί την κατανάλωση, δεδομένων των αποφάσεων της προηγούμενης ημέρας, ακολουθώντας το μοντέλο 
οικονομικής κατανομής.
Το μοντέλο δέσμευσης μονάδων που περιγράφεται παρακάτω, θα χρησιμοποιηθεί για την κατασκευή των 
προσφορών που θα αποτελούν την είσοδο στο μοντέλο του αλγορίθμου EUPHEMIA που θα ακολουθήσει. 
*στρατηγική*

\section{Μοντέλο Βελτιστοποίησης Δέσμευσης Μονάδων}

Το πρόβλημα δέσμευσης μονάδων για ένα σύνολο γεννητριών G ορίζεται σε ένα χρονικό ορίζοντα προγραμματισμού T χρονικών περιόδων, κατά τις οποίες εισαγάγουμε δυαδικές μεταβλητές δέσμευσης ugt. Αν μια μονάδα παραγωγής g είναι εντός λειτουργίας τη χρονική περίοδο t, τότε ugt=1, διαφορετικά, αν είναι εκτός λειτουργίας, ugt= 0. Η ποσότητα παραγόμενης ισχύος pgt και η ποσότητα εφεδρείας rgt της μονάδας τη χρονική περίοδο t αποτελούν συνεχείς μεταβλητές απόφασης. 
*TODO: Οι εφεδρείες συγκεκριμένα θα παραλειφθούν στην υλοποίησή των μοντέλων Δέσμευσης Μονάδων και EUPHEMIA, καθώς στην ευρωπαϊκή αγορά ενέργειας δε λαμβάνονται υπόψη σε αυτό το στάδιο προγραμματισμού. *s

\section{Αρχικές Συνθήκες}

Καθώς η δέσμευση μονάδων αποτελεί διαρκή διαδικασία, για την επίλυσή ενός προβλήματος σε χρονικό ορίζοντα T απαιτούνται δεδομένες συνοριακές συνθήκες για τις μονάδες. Συγκεκριμένα, για έναν ορίζοντα {1,...T} απαιτείται ένας πρότερος ορίζοντας To, η διάρκεια του οποίου ισοδυναμεί με το μέγιστο χρόνο επίδρασης μιας απόφασης σε μελλοντική απόφαση. Για παράδειγμα, εάν ο χρόνος διάρκειας παραμονής μιας μονάδας εκτος λειτουργίας (minimum downtime) από τη στιγμή που τίθεται εκτός είναι 16 ώρες, τότε το Toθα πρέπει να έχει εύρος τουλάχιστον 16 ώρες. Συμβολίζουμε τις αρχικές αποφάσεις δεσμευσης και παραγωγής ως u0g, p0g, αντίστοιχα, για δεδομένη γεννήτρια gEG.

\section{Μεταβάσεις}

Για την περιγραφή των περιορισμών προγραμματισμού, είναι απαραίτητος ο ορισμός μεταβλητών περιγραφής της εκκίνησης και διακοπής λειτουργίας μιας γεννήτριας gEG. Με vgt συμβολίζουμε τη δυϊκή μεταβλητή απόφασης εκκίνησης μιας γεννήτριας g τη χρονική περίοδο t, ενώ με zgt συμβολίζουμε τη δυϊκή μεταβλητή τερματισμού λειτουργίας μιας γεννήτριας g τη χρονική περίοδο t. Οι εν λόγω μεταβλητές συσχετίζονται με τις μεταβλητές δέσμευσης μονάδας ως εξής:

ug,t = ug,t-1+ vgt - zgt gEG, t = 2, ...T

Δηλαδή, μία γεννήτρια μπορεί
-  να μεταβεί από κατάσταση δέσμευσης σε κατάσταση διακοπής, άρα
-ugt=0
-ugt-1=1
-vgt=0
- zgt=1
- να μεταβεί από κατάσταση μη δέσμευσης σε κατάσταση δέσμευσης, όπου
	-ugt=1
	-ugt-1=0
	-vgt=1
	-zgt=0
- να παραμείνει στην ίδια κατάσταση με την προηγούμενη περίοδο
	ugt=ugt-1

Για την τρίτη περίπτωση, παρατηρούμε πως ικανοποιείται η εξίσωση για vgt=zgt=1. Για να αποφύγουμε τη αυτή την περίπτωση ταυτόχρονης εκκίνησης και τερματισμού της γεννήτριας, προσθέτουμε τον περιορισμό vgt+zgt<=1.

\section{Minimum Uptime και Downtime}

Μετά την ενεργοποίηση μιας γεννήτριας, συνήθως πρέπει να παρέλθει ένα ελάχιστο χρονικό διάστημα ώστε να είναι δυνατή η απενεργοποίησή της, λόγω τεχνικών περιορισμών, κάλυψης του κόστους εκκίνησης, και διατήρησης της λειτουργικής τους απόδοσης. Αντίστοιχα, μία γεννήτρια απαιτεί έναν ελάχιστο χρόνο για να ψυχρανθεί από την ώρα που τέθηκε εκτός λειτουργίας.

O ENTSO-e, ο Ευρωπαϊκός Σύνδεσμος Διαχειριστών Συστημάτων Μεταφοράς Ηλεκτρικής Ενέργειας, είναι ο σύνδεσμος συνεργασίας των ευρωπαϊκών διαχειριστών συστημάτων μεταφοράς ηλεκτρικής ενέργειας (TSOs). Τα 40 μέλη TSOs, που εκπροσωπούν 36 χώρες, είναι υπεύθυνα για την ασφαλή και συντονισμένη λειτουργία του συστήματος ηλεκτρικής ενέργειας της Ευρώπης, το οποίο αποτελεί το μεγαλύτερο διασυνδεδεμένο ηλεκτρικό δίκτυο στον κόσμο. Εκτός από τον βασικό και ιστορικό ρόλο του στην τεχνική συνεργασία, o ENTSO-E αποτελεί επίσης τη συλλογική φωνή των TSOs.


/σσtion{Δομή της εργασίας}



\chapter{Θεωρητικό υπόβαθρο}

Ο καπετάν Φουρτούνας, ένας από τους πέντε δημογέροντες του χωριού, πρόβαλε
στην πόρτα του καφενέ. Αψηλός, κορμάτος, παλιός καραβοκύρης, που χρόνια
αλώνιζε τη Μαύρη Θάλασσα κουβαλώντας ρούσικο σιτάρι και κάνοντας
κοντραμπάντο. Τρίχα δεν είχε το πρόσωπο του· σπανός και μαυροκίτρινος,
ταγαριασμένος, με βαθιές ζαρωματιές, και τα μάτια του τα μικρά, τα
κατάμαυρα, σπίθιζαν. Γέρασε, γέρασε και το καράβι του, τσακίστηκε μια νύχτα
απόξω από την Τραπεζούντα, κι ο καπετάν Φουρτούνας, καραβοτσακισμένος,
μπουχτισμένος, γύρισε στο χωριό του να σουρώσει όσο μπορεί περισσότερη ρακή
και, σαν έρθει η ώρα η καλή, να γυρίσει το πρόσωπό του κατά τον τοίχο και να
πεθάνει. Πολλά είχαν δει τα μάτια του, βαρέθηκε· δε βαρέθηκε, κουράστηκε, μα
ντρέπουνταν να το μολοήσει.


\section{Σημασιολογία γλωσσών προγραμματισμού}

Φορούσε σήμερα τις αψηλές καπετανίστικες μπότες του, τον κίτρινό του μουσαμά
και το αρχοντικό καλπάκι από αληθινό αστρακάν. Κρατούσε και το αψηλό ραβδί
του του δημογέροντα. Δυο τρεις χωριανοί προσηκώθηκαν να τον καλέσουν να
πάρει μια ρακή.

--- Δεν έχω καιρό, παιδιά, μήτε για ρακή, είπε· Χριστός ανέστη! Πάω στο
αρχοντικό του παπά, όπου έχουμε σύναξη οι προεστοί. Σε μια ώρα να κοπιάσετε
κι όσοι από σας είστε καλεσμένοι· κάντε το σταυρό σας κι ελατέ, καλά το
κατέχετε, έχουμε σήμερα δουλειά. Κι ένας να πάει να φωνάξει τον Παναγιώταρο
το σαμαρά, με του διαόλου τα γένια· αυτόν τον έχουμε μεγάλη ανάγκη.

Σώπασε μια στιγμή, τα μάτια του έπαιξαν παμπόνηρα:

--- Αν δεν είναι σπίτι του, θα~’ναι στης χήρας, είπε κι όλοι ξέσπασαν στα
γέλια.

Μα ο γερο-Χριστοφής ο αγωγιάτης, που~’χε μάθει στα νιάτα του, κι ας το
πλέρωσε ακριβά, τι θα πει σεβντάς, πετάχτηκε απάνω:

--- Τι γελάτε, μωρέ σερσέμηδες; φώναξε. Καλά κάνει· φωτιά στα τόπια, μωρέ
Παναγιώταρε, και μην τους ακούς! Λίγη είναι η ζωή, πολύς ο θάνατος, βίρα!

Ο χοντρο-Δημητρός ο χασάπης κούνησε το φρεσκοξουρισμένο κεφάλι:

--- Ο Θεός να~’χει καλά τη χήρα την Κατερίνα μας, είπε. Ο διάολος ξέρει από
τι κέρατα μας γλιτώνει!

Ο καπετάν Φουρτούνας γέλασε.

--- Βρε παιδιά, είπε, μη μαλώνετε. Χρειάζεται μια παστρικιά σε κάθε χωριό,
να μη βρίσκουν τον μπελά τους οι τίμιες. Είναι σαν τη βρύση του δρόμου,
μαθές· περνούν και πίνουν οι διψασμένοι· αλλιώς θα χτυπούσαν αράδα τις
πόρτες μας. Κι οι γυναίκες, όταν τους ζητήσεις νερό...

Στράφηκε, είδε το δάσκαλο:

--- Χατζη-Νικολή μου, του κάνει, ακόμα είσαι εδώ; Προεστός είσαι και του
λόγου σου, έχουμε σύναξη. Σκολειό έκαμες και τον καφενέ, σκόλασε κι έλα!

--- Να~’ρθω κι εγώ; είπε ο γερο-Χριστοφής κι έπαιξε το μάτι στην παρέα· κάνω
για Ιούδας.

Μα ο καπετάν Φουρτούνας είχε πάρει την ανηφοριά, βαριακουμπώντας το ραβδί
του στο καλντιρίμι. Δεν ήταν στα καλά του σήμερα· τον έσφαζαν πάλι οι
ρεματισμοί, μάτι δεν είχε κλείσει όλη νύχτα. Ήπιε πρωί πρωί δυο τρία
νεροπότηρα ρακή, για γιατρικό, μα του κάκου· οι πόνοι δε σκολνούσαν. Μήτε η
ρακή τους έβανε κάτω.

--- Μωρέ, αν δε ντρέπουμουν, μουρμούρισε, κι άρχιζα τις φωνές, μπορεί
ν’~αλάφρωναν οι πόνοι· μα έλα που δε με αφήνει το παντέρμο το φιλότιμο! Και
πρέπει να περπατώ ντούρος και να κάνω πως γελώ. Κι αν πέσει χάμω το ραβδί
μου, να μην αφήσω κανέναν κερατά να με βοηθήσει, μα να σκύψω μοναχός μου να
το πιάσω. Δάγκανε, μωρέ καπετάν Φουρτούνα, τα χείλια, όρτσα τα πανιά,
κατακέφαλα στο κύμα, βάρδα μην ντροπιαστείς! Μπόρα είναι, μαθές, κι η ζωή,
θα περάσει!

Έγρουζε και σιγοβλαστημούσε σκαμπανεβαίνοντας. Στάθηκε μια στιγμή, κοίταξε
γύρα, κανένας δεν τον έβλεπε· αναστέναξε κι αλάφρωσε λίγο. Κοίταξε κατά
πάνω, είδε στην κορυφή του χωριού ν’~ασπρογυαλίζει ανάμεσα από τα δέντρα το
σπίτι του παπά με τα λουλακιά παραθυρόφυλλα.

--- Πήγε κι έχτισε στην κορυφή του χωριού ο διαολόπαπας! μουρμούρισε· την
κατάρα του να~’χω!

Και πήρε πάλι τον ανήφορο.

Στο σπίτι του παπά είχαν κιόλα φτάσει δυο από τους προεστούς και κάθουνταν
σταυροπόδι στο μεντέρι, αμίλητοι, και περίμεναν τα τραταρίσματα. Ο παπάς
είχε μπει στην κουζίνα να δώσει διαταγές, κι η μοναχοκόρη του η Μαριορή
ετοίμαζε το δίσκο με τον καφέ, το δροσερό νερό και το γλυκό του κουταλιού.

Πλάι στο παραθύρι είχε θρονιαστεί ο πρώτος δημογέροντας της Λυκόβρυσης, από
μεγάλο τζάκι, αρχοντάνθρωπος, καλοθρεμμένος, με τσόχινο σαλβάρι, με
χρυσοκέντητο μεϊντάνι και μ’~ένα χοντρό χρυσό δαχτυλίδι στο δείχτη του
χεριού του --- η βούλα του με τα δυο κεφαλαία σφιχταγκαλιασμένα γράμματα:
Γ. Π., Γεώργιος Πατριαρχέας. Τα χέρια του ήταν παχιά και μαλακά, σα
Δεσπότη. Δε δούλεψε ποτέ του, είχε ένα τσούρμο φαμέγιους και κολίγους που
δούλευαν και τον τάιζαν. Και χόντρυνε τ’~άντερό του, φάρδυναν τα καπούλια
του, είχαν γίνει σαν της φοράδας, κρεμάστηκαν οι κοιλιές του, και τα
προγούλια του κατέβαιναν τρεις πατωσιές κι αναπαύουνταν, το ένα απάνω στο
άλλο, στο μαλλιαρό αφράτο απανωστήθι. Δυο τρία δόντια μπροστινά του~’λειπαν,
άλλο κουσούρι δεν είχε, κι όταν μιλούσε, φαφλάτιζε και μπερδεύουνταν· μα και
το κουσούρι αυτό πλήθαινε την αρχοντιά του, γιατί σε ανάγκαζε, όταν μιλούσε,
να σκύβεις για να καταλάβεις τι λέει.

Δεξά του, στη γωνιά, αδύναμος, λιγδοτάμπαρος, φτενοκέφαλος, με τσιμπλιασμένα
μάτια, με δυο χοντρές χερούκλες γεμάτες ρόζους, κάθουνταν συμμαζεμένος και
ταπεινός ο δεύτερος προεστός, ο πιο βαρβάτος νοικοκύρης του χωριού, ο
γερο-Λαδάς. Εβδομήντα χρόνια σκύβει στη γης, τη σκάβει, τη σπέρνει, τη
θερίζει, της φυτεύει ελιές κι αμπέλια, τη στύβει και της πίνει το αίμα. Ποτέ
του, από μικρό κουτσούβελο, δεν ξεκόλλησε από τη γης. Αχόρταγος, λιμασμένος,
έπεφτε απάνω της, της έδινε ένα και της ζητούσε χίλια, και ποτέ του δεν
έλεγε: «Δόξα σοι ο Θεός», μα πάντα μουρμούριζε, αφχαρίστητος. Και στα
γερατιά του δεν τον έφτανε πια η γης· όσο ζύγωνε στο θάνατο κι ένιωθε πως
λίγα πια ήταν τα καρβέλια του, βιάζουνταν να προλάβει και να φάει τον
κόσμο. Κίνησε λοιπόν να δανείζει με βαρύ διάφορο τους χωριανούς του· έβανα
οι δύστυχοι αμανάτι τ’~αμπελοχώραφά τους και τα σπίτια, έρχουνταν ο καιρός
της πληρωμής, δεν είχαν να πλερώσουν, έβγαιναν τα πράματα στο σφυρί, και
τα’~χαφτε ο γερο-Λαδάς.

Κι όλο κλαίγουνταν, κι όλο πεινούσε, κι η γυναίκα του γύριζε ξυπόλυτη, και
μια θυγατέρα, που~’χε καταφέρει να κάμει, την άφησε κι αυτή να πεθάνει,
γιατί δεν κάλεσε, όταν έπεσε στο κρεβάτι, γιατρό να τη δει.

«Πολλά έξοδα, είπε, οι μεγάλες πολιτείες είναι μακριά, πού να φέρνεις
γιατρό! Κι ύστερα, τι ξέρουν κι αυτοί; τον κακό τους τον καιρό! Έχουμε εδώ
τον παπά μας, αυτός κατέχει παλιά γιατροσόφια, θα τον πλερώσω να κάμει κι
ευκέλαιο, και θα γιάνει, και θα κοστίσει και πιο φτηνά.»

Μα τα ματζούνια του παπά πήγαν χαμένα, το ευκέλαιο δεν έπιασε, κι η κοπέλα
πέθανε, δεκαφτά χρονών, και γλίτωσε από τον κύρη της· γλίτωσε κι αυτός από
τα πολλά τα έξοδα του γάμου. Μια μέρα, λίγους μήνες μετά το θάνατό της,
κάθισε και τα λογάριασε: Προίκα τόσο απάνω κάτω, ρουχισμός, τραπέζια,
καρέκλες, τόσα, θ’~αναγκάζουνταν να καλέσει στο γάμο τους συγγενείς, κι
αυτοί θένε να φάνε τον περίδρομο, βάλε κρέατα, ψωμιά, κρασιά, τόσα... Έκαμε
τη σούμα, πάρα πολλά έξοδα, θα τον ξετίναζε η θυγατέρα του, δεν πειράζει το
λοιπόν, όλοι θα πεθάνουμε... Γλίτωσε κι από τα βάσανα του κόσμου – άντρες,
παιδιά, αρρώστιες, μπουγάδες. Τυχερή στάθηκε, ο Θεός σχωρέσ’~τη!


\section{Θεωρία πεδίων}

Μπήκε η Μαριορή με το δίσκο, χαιρέτησε χαμοβλεπούσα τους προεστούς, στάθηκε
μπροστά από τον άρχοντα. Χλωμή, μεγαλομάτα, γαϊτανοφρύδα, με δυο χάντρες
πλεξούδες καστανά μαλλιά, τυλιμένες στεφανωτά γύρα στο κεφάλι. Γέμισε ο
γερο-άρχοντας τρουλωτά το κουταλάκι του γλυκό βύσσινο, κοίταξε την κοπέλα,
σήκωσε το ποτήρι.

--- Στις χαρές σου, Μαριορή μου, ευκήθηκε· ο γιος μου βιάζεται.

Ήταν αρραβωνιασμένη η παπαδοπούλα με τον μοναχογιό του το Μιχελή, και ο
παπάς καμάρωνε που θα’~κανε τέτοιο συμπεθεριό και θα’~πιανε γρήγορα αγγόνια.

--- Δεν μπορώ να καταλάβω γιατί βιάζεται τόσο ο βλογημένος· δεν μπορεί,
λέει, πια... πρόσθεσε γελώντας ο άρχοντας κι έκλεισε το μάτι στην κοπέλα.

Κι αυτή κοκκίνησε ως το λαιμό, πλαντούσε, δεν μπόρεσε να μιλήσει.

--- Χαρές να’~χουμε! είπε ο παπα-Γρηγόρης, μπαίνοντας με μια μπουκάλα
μοσχάτο κρασί. Με την ευκή του Χριστού και της Παναγιάς!

Άγριος, κοτσονάτος, με διχαλωτή γενειάδα κάτασπρη, καλοθρεμμένος, μύριζε
λιβάνι και βούτυρο. Είδε την κόρη του να κοκκινίζει και, για ν’~αλλάξει
κουβέντα:

--- Πότε με το καλό θα παντρέψεις και την παρακόρη σου το Λενιό; ρώτησε.

Το Λενιό ήταν μια από τις μπασταρδοπούλες που ο άρχοντας είχε σκαρώσει με
τις παραδουλεύτρες του. Την είχε αρραβωνιάσει με τον ήμερο πιστό τσοπάνη
του, το Μανολιό, και την είχε αρχοντικά προικίσει μ’~ένα κοπάδι πρόβατα, που
τα’~βοσκε ο Μανολιός στο αντικρινό βουνό της Παναγιάς.

--- Αν θέλει ο Θεός, αποκρίθηκε, αυτές τις μέρες· το Λενιό, λέει,
βιάζεται. Βιάζεται το καλορίζικο· σηκώθηκε, μαθές, το βυζί του και θέλει να
βυζάξει γιο. «Μάης μπαίνει, μου’~πε προχτές, Μάης μπαίνει, αφεντικό, και
πρέπει να βιαστούμε.»

Γέλασε πάλι καλόκαρδα, τα προγούλια του σείστηκαν.

--- Το Μάη είπε, παντρεύονται τα γαϊδούρια, έχει δίκιο το Λενιό· πρέπει να
βιαστούμε. Άνθρωποι είναι κι αυτοί, και ας είναι και φαμέγοι.

--- Καλός είναι ο Μανολιός, έκαμε ο παπάς· καλά θα ζήσουν.


%\englishtext

\selectlanguage{greek}


%%%  Bibliography

% You shouldn't want to include all the contents of thesis.bib
% in your bibliography (do you?)
\nocite{*}

\bibliographystyle{softlab-thesis}
\bibliography{thesis}


%%%  Appendices

\backmatter

\appendix

\chapter{Ευρετήριο συμβολισμών}

$A \rightarrow B$ : συνάρτηση από το πεδίο $A$ στο πεδίο $B$.

\chapter{Ευρετήριο γλωσσών}

\begin{description}
\item[C++:] πώς θα βγάλω λεφτά...
\item[Haskell:] η γλώσσα της ζωής μου αλλά πάνε οι μπύρες...
\item[Javascript:] χα, χα, χα...
\item[Python:] πώς θα τελειώνω για να πάω για μπύρες...
\end{description}


\chapter{Ευρετήριο αριθμών}

\begin{description}
\item[17:] ask Zachos.
\item[42:] life, the universe and everything --- ask Douglas.
\end{description}


%%%  End of document

\end{document}
